\section{Simplicial recognition and the two topologies}
\label{sec:simplicial-comparison}

\subsection{Connected boundaries}
For any poset $P$ and any $x\in P$, put $B_x=P_{<x}$.
The following elementary equivalence specifies the connectedness
condition that is weaker than the spherical hypothesis in
Theorem~\ref{thm:cw}.

\begin{proposition}\label{prop:boundary-connected}
The following are equivalent: $B_x$ is nonempty and connected
in either Alexandrov topology; its undirected comparability graph
is nonempty and connected; and $|\Delta(B_x)|$ is nonempty and
connected. If the principal ideal of $x$ is finite, the
comparability graph can be replaced by the Hasse graph.
\end{proposition}
\begin{proof}
A comparability component is both an upper and a lower set, hence
is clopen in either Alexandrov topology. Conversely a clopen
separation cannot separate comparable points: lower openness
of the two sides, or dually upper openness, rules this out.
The comparability edges are exactly the edges of the order
complex. Their graph components partition its simplices and give
clopen subcomplexes in the weak topology. A finite edge path is
a topological path, proving the second equivalence. In a finite
principal ideal, comparable points can be joined by a saturated
chain. The nonempty qualification excludes the empty boundary of
a vertex; the boundary of a regular edge is the disconnected
space $S^0$.
\end{proof}
Thus a connected-boundary axiom in dimensions at least two means
that any two proper faces are joined by a finite incidence zigzag
among proper faces. It is not a replacement for sphericality.

\subsection{The global intersection condition for simplices}
We prove the first assertion of
Theorem~\ref{thm:simplicial-comparison}. For the nonempty face
poset of an abstract simplicial complex, adjoining the empty face
gives binary meets by intersection and gives a finite Boolean
interval below each face.

Conversely assume the finite Boolean interval condition, and
write $V_x$ for the atoms of $P^0$ below $x$. If binary meets exist
and $V_x=V_y$, all these atoms lie below $x\wedge y$.
Within either Boolean interval they generate its top, so
$x=x\wedge y=y$. Thus the atom-set map is injective.
If instead injectivity is assumed, every subset of $V_x$ is
the atom set of an element of $[0,x]$. Hence the nonempty atom
sets form an abstract simplicial complex. Moreover
$V_x\subseteq V_y$ implies $x\le y$: the element below $y$
with atom set $V_x$ must be $x$. The atom-set map is an order
isomorphism to the nonempty face poset of that complex.
Intersections supply its binary meets. This proves all three
equivalent formulations, including the empty case.

Local Booleanity alone is insufficient. Two distinct edges with
the same two endpoints form a regular circle with Boolean
principal augmented intervals, but the two edge cells have the
same vertex set. Their common lower faces have no greatest
element, so adjoining zero does not give a meet-semilattice.

\begin{proposition}[A weaker predecessor condition]\label{prop:predecessor-repair}
The conjunction of binary meets in $P^0$ and the condition that
every nonminimal cell has at least two maximal proper faces does
not characterize simplicial complexes, even among finite regular
CW face posets. Replacing the predecessor condition by finite
Boolean principal intervals gives the exact criterion.
Equivalently, one may require finite distributive principal
intervals together with that predecessor condition.
\end{proposition}
\begin{proof}
Every positive-dimensional regular cell satisfies the predecessor
condition: a unique maximal proper face would make its boundary
order complex a cone, whereas that complex is a nonempty sphere.
The face poset of a square has binary intersections of closed
faces and satisfies the condition. Its augmented top interval
has ten elements and four atoms, so it cannot be the sixteen-element
Boolean lattice on those atoms. Even regularity does not follow
from the two weaker conditions: three incomparable vertices below
a single top cell satisfy them, but that cell has three boundary
points rather than the two of a regular one-cell.

In a finite distributive lattice every nonzero join-irreducible
has exactly one lower cover. The predecessor condition therefore
makes every such element an atom. Every element is a join of
join-irreducibles, so the lattice is atomistic. Distinct atoms have
zero meet; distributivity shows that different subsets of atoms
have different joins. Every subset has a join and every element
is such a join. This gives a Boolean lattice. The converse is
immediate, proving the final equivalence.
\end{proof}
In a regular cell complex, the meet condition says that two closed
cells intersect in a closed cell or not at all. It makes no claim
that the cells can be realized by convex polytopes.

\subsection{Flags and interval incidences}
The finite nonempty chains of a poset $L$, ordered by inclusion,
are exactly the nonempty faces of $\Delta(L)$. Subchains are
chains, and intersection gives the augmented meets. This is a
simplicial construction in every dimension.

The point-and-interval construction has a different incidence
relation. Its vertex cells are $P_a$ and its one-cells are
$I_{ab}$ for $a<b$, with the two endpoint incidences
$P_a,P_b<_{\mathrm{face}}I_{ab}$. It is the one-dimensional
simplicial complex of the comparability graph of $L$. It has
no higher-dimensional cells even when $L$ has long chains.
Its positional order, used in Theorem~\ref{thm:flags}, instead
puts $P_a<_{\mathrm p}I_{ab}<_{\mathrm p}P_b$.

A finite \emph{prefigure} in this construction is a positional
chain; a \emph{figure} is one closed under endpoint incidence.
The endpoint-closure operation is
\[
 F\longmapsto F\cup\bigcup_{I_{ab}\in F}\{P_a,P_b\}.
\]
A symbol preceding $I_{ab}$ has its last endpoint at most $a$,
and a symbol following it has its first endpoint at least $b$.
Thus endpoint closure preserves the positional-chain condition.
For example, $\{I_{ab}\}$ is a prefigure but is not a figure.
This also explains why the finite admissibility theorem was
stated for all pairwise orthogonal subsets, without assuming
incidence closure.

\subsection{The weak comparison with an Alexandrov space}
We finish Theorem~\ref{thm:simplicial-comparison} by checking the
local criterion in McCord's Theorem~6 \cite[p.~467]{McCord}.
That theorem says that a map is a weak equivalence if its
restrictions over a basis-like open cover are weak equivalences.
For $U_x=L_{\le x}$, the inverse image under $\mu_L$ consists
of points with a positive coordinate at some vertex of $U_x$.
It is open, being the complement of the full subcomplex on
$L\setminus U_x$. Renormalize the coordinates in $U_x$ and
linearly suppress the others. This deforms the inverse image
onto $|\Delta(U_x)|$, a cone with apex $x$. The same greatest
element contracts the Alexandrov space $U_x$. The principal
lower sets form a basis, so McCord's local theorem applies.
The coordinate deformation is continuous simplex by simplex
in the weak topology; no local finiteness is required.

This is generally only a weak equivalence. Take two incomparable
minima below two incomparable maxima. Its order complex is a
circle. Every continuous map from that finite connected poset
to a Hausdorff space is constant: comparable points must have
equal images. A homotopy inverse into the circle would therefore
make its identity null-homotopic, which is impossible.
For a regular CW face poset, the order complex instead gives
a barycentric subdivision homeomorphic to the original
Hausdorff realization. These comparisons have different targets.
